\cleardoublepage
\phantomsection
\addcontentsline{toc}{chapter}{Liste des figures}
\listoffigures

\cleardoublepage
\phantomsection
\addcontentsline{toc}{chapter}{Liste des tableaux}
\listoftables

\chapter*{Résumé}
\addcontentsline{toc}{chapter}{Résumé}

La robotique en essaim cherche à obtenir des comportements collectifs
à partir de règles locales appliquées par de nombreux robots. Former une
forme représente un défi particulier, puisque chaque robot ne dispose que
d’informations limitées et doit pourtant contribuer à un objectif commun.
Les méthodes existantes reposent soit sur des règles écrites à la main,
soit sur des contrôleurs appris, et il reste difficile de les comparer,
parce que leurs hypothèses et leurs protocoles d’évaluation diffèrent.

Ce mémoire présente un contrôleur décentralisé fondé sur un réseau de
neurones sur graphes (GNN), partagé par tous les robots. Chaque robot
s’appuie sur son état local, sur les informations de ses voisins et sur
une consigne commune qui décrit la forme à réaliser. Nous avons mis en
place un simulateur bidimensionnel ainsi qu’un protocole d’évaluation qui
teste la formation de formes, le retrait de robots et les variations de
taille de l’essaim.

Des essais préliminaires menés avec un prototype antérieur suivent mieux
un contour circulaire qu’une règle locale de regroupement. Ces
observations ont toutefois une portée limitée, car les deux approches ne
reçoivent pas la même information. L’entraînement et l’évaluation
comparative de l’implémentation actuelle restent encore à réaliser.

Ce travail propose un cadre expérimental pour étudier le rôle que jouent
les GNN dans le contrôle décentralisé des essaims de robots, accompagné
d'un protocole conçu pour que les expériences soient faciles à reproduire.

\noindent\textbf{Mots-clés :} robotique en essaim, contrôle décentralisé,
réseaux de neurones sur graphes, formation de formes, apprentissage
différentiable, simulation robotique.

\chapter*{Abstract}
\addcontentsline{toc}{chapter}{Abstract}

Swarm robotics tries to get collective behaviour out of simple local rules, applied by many robots at once. Shape formation is a tricky case. Each robot sees only a small part of the picture, yet all of them have to work toward the same goal. Existing approaches rely either on hand-written rules or on learned controllers, and it is hard to compare the two, since their assumptions and evaluation methods rarely match.

This thesis looks at a decentralized controller built on a graph neural network (GNN), shared by every robot. Each robot uses its own state, what its neighbours tell it, and a common instruction that describes the target shape. We test it in a two-dimensional simulator, with a protocol covering shape formation, robot removal and changes in swarm size.

Early tests with a previous prototype suggest that it follows a circular contour better than a local gathering rule. That said, these results should be read with caution. The two approaches do not get the same information about the target. Training and comparative evaluation of the current implementation have not been done yet.

What this work offers, in the end, is an experimental framework for asking what GNNs actually bring to decentralized swarm control. The protocol is designed so that others can reproduce the experiments without much trouble.

\noindent\textbf{Mots-clés.} robotique en essaim, contrôle décentralisé,
réseaux de neurones sur graphes, formation de formes, simulation robotique.

\cleardoublepage
\tableofcontents